\subsection{Por qué sockets fue más rápido}

H1 salió como esperábamos, y con una diferencia más grande de lo que pensábamos: con 100~bytes y un cliente, gRPC tardó unas 6,5 veces más que sockets en la mediana. En \emph{localhost} el paquete casi no viaja, así que esos 0,25~ms de más por mensaje salen del software. Con sockets, cada envío es un \texttt{json.dumps}, una línea por TCP y un \texttt{json.loads} del acuse. Con gRPC, cada llamada pasa además por las tramas y los encabezados de HTTP/2, los metadatos, el plazo de la llamada y el hilo del \texttt{ThreadPoolExecutor} que la atiende en el servidor. No medimos cada una de esas capas por separado, así que no podemos decir cuál pesa más. Lo que sí vemos es que el costo es fijo por llamada, porque se repite casi igual en las cinco repeticiones (de 0,288 a 0,295~ms).

Esto coincide con lo que reportan Zhu \emph{et al.}~\cite{Zhu2023}: en su medición, interpretar el protocolo de gRPC costaba bastante más que tratar el mismo tráfico como TCP. Y no contradice a Niswar \emph{et al.}~\cite{Niswar2024}, donde gRPC fue más rápido que REST y GraphQL, porque en ese estudio las otras opciones también pasaban por HTTP. Aquí sockets no tiene ninguna capa encima de TCP.

\subsection{Efecto del tamaño del mensaje}

H2 se cumplió, pero por un motivo distinto del que planteamos. Esperábamos que la diferencia en milisegundos se mantuviera y que pesara menos porque el total crecía en los dos mecanismos. La diferencia sí se mantuvo (de 0,25 a 0,23~ms), pero lo que bajó el cociente fue que sockets se hizo más lento con 10~KB (un 60\,\% más) y gRPC casi no cambió (un 4\,\% más). Creemos que la causa está en dónde se procesa el mensaje. En sockets, los 10~KB pasan por \texttt{json.dumps}, por la lectura de la línea y por \texttt{json.loads}, y parte de ese trabajo se hace en Python. En gRPC la serialización la hace la biblioteca de Protocol Buffers, cuya implementación para Python corre en código compilado. No lo comprobamos con un perfilador, así que lo dejamos como explicación probable. Lo que sí sale de los datos es que, en este rango de tamaños, el costo fijo de cada llamada de gRPC pesa mucho más que el costo de copiar 10~KB.

\subsection{Varios clientes a la vez}

H3 no se cumplió. Habíamos supuesto que con tres clientes gRPC sufriría más por el GIL, porque cada llamada pasa por más código Python de la biblioteca. Los datos no muestran eso: con 100~bytes, la mediana de sockets subió un 24\,\% y la de gRPC un 26\,\%, prácticamente lo mismo. Una posible explicación es que en gRPC la red y HTTP/2 los maneja el núcleo en C, que no necesita el GIL, y eso compensa el código Python extra. Por otro lado, la máquina tiene 28 hilos de procesador y los tres clientes son procesos separados, así que solo los servidores comparten el GIL.

Donde gRPC sí se acercó a sockets fue en los mensajes por segundo. Con tres clientes, gRPC atendió algo más del doble que con uno (de 2640 a 5636), y sockets un 59\,\% más (de 7747 a 12\,306). Aun así, sockets siguió atendiendo más mensajes por segundo en todas las condiciones. Este resultado hay que leerlo con cuidado, porque los mensajes por segundo de sockets varían mucho entre repeticiones.

\subsection{Los envíos lentos de sockets}

El resultado que no esperábamos es que sockets, con la mediana más baja, tenga los envíos más lentos de todo el experimento: 13 de 4000 pasaron de 10~ms y uno llegó a 27,6~ms, mientras que con gRPC ninguno pasó de 7,8~ms. Aparecen en 9 corridas distintas y en las cuatro condiciones, así que no vienen de una sola corrida que haya salido mal. No encontramos la causa; puede ser la planificación de hilos de Windows o algo del servidor de sockets, pero no lo medimos. Son el 0,3\,\% de los envíos, así que casi no mueven la mediana ni el percentil~95, pero sí la media y los mensajes por segundo. Por eso el criterio de decisión usa medianas y rangos por repetición, y no la media de todos los envíos. Si lo que importa es una latencia predecible, gRPC fue más estable en este experimento.

\subsection{Convergencia de los relojes de Lamport}
\label{sec:discusion-lamport}

El criterio de H4 decía que, si los intercambios no daban 2, 4 y 6, había que revisar la ronda antes de discutir la hipótesis. La revisamos. La hipótesis describe una ronda con una sola secuencia de mensajes: el primer nodo envía al segundo, este al tercero y así hasta volver al primero. Pero la ronda que quedó en el código el 20 de noviembre, antes de la corrida oficial, usa dos fichas que salen al mismo tiempo de dos nodos distintos. Así los dos primeros envíos son concurrentes, que es el caso que pide el enunciado. El camino que la información tiene que recorrer para llegar a todos es el mismo, pero mientras tanto la otra ficha también envía mensajes, y esos mensajes entran en la cuenta. Por eso salen 6 y 8 intercambios en vez de 4 y 6. Con dos nodos no cambia nada, porque basta con que cada uno reciba un mensaje del otro. No cambiamos la hipótesis después de ver los datos: la reportamos como estaba escrita y dejamos aquí la diferencia.

Lo que sí se mantiene es la tendencia. Los intercambios crecen de 2 a 6 y a 8, y siempre dan lo mismo en las cinco repeticiones, porque el patrón de envíos es fijo.

En el tiempo, H4 no se pudo confirmar por las 8 corridas que tardaron más de medio segundo. Revisamos una de ellas (sockets, 2 nodos, repetición~5). El nodo~A registró su primer envío y la recepción correspondiente en el otro nodo llegó 0,51~s después. Lo que pasó es que el nodo~A intentó conectarse antes de que el nodo~B abriera su puerto, y en Windows ese intento fallido tarda alrededor de medio segundo antes de que el nodo vuelva a intentar. Ese tiempo es arranque de los procesos y no convergencia de los relojes. Con gRPC, donde solo una corrida tuvo este problema, la tendencia sí se ve: las medianas son 5,7, 11,2 y 13,0~ms para 2, 3 y 4 nodos. Una corrección sería que los nodos esperen a que todos los servidores estén escuchando antes del primer envío. No repetimos la corrida con ese cambio, y lo dejamos como amenaza a la validez.
